% Options for packages loaded elsewhere
\PassOptionsToPackage{unicode}{hyperref}
\PassOptionsToPackage{hyphens}{url}
\documentclass[
]{article}
\usepackage{xcolor}
\usepackage[margin=1in]{geometry}
\usepackage{amsmath,amssymb}
\setcounter{secnumdepth}{-\maxdimen} % remove section numbering
\usepackage{iftex}
\ifPDFTeX
  \usepackage[T1]{fontenc}
  \usepackage[utf8]{inputenc}
  \usepackage{textcomp} % provide euro and other symbols
\else % if luatex or xetex
  \usepackage{unicode-math} % this also loads fontspec
  \defaultfontfeatures{Scale=MatchLowercase}
  \defaultfontfeatures[\rmfamily]{Ligatures=TeX,Scale=1}
\fi
\usepackage{lmodern}
\ifPDFTeX\else
  % xetex/luatex font selection
\fi
% Use upquote if available, for straight quotes in verbatim environments
\IfFileExists{upquote.sty}{\usepackage{upquote}}{}
\IfFileExists{microtype.sty}{% use microtype if available
  \usepackage[]{microtype}
  \UseMicrotypeSet[protrusion]{basicmath} % disable protrusion for tt fonts
}{}
\makeatletter
\@ifundefined{KOMAClassName}{% if non-KOMA class
  \IfFileExists{parskip.sty}{%
    \usepackage{parskip}
  }{% else
    \setlength{\parindent}{0pt}
    \setlength{\parskip}{6pt plus 2pt minus 1pt}}
}{% if KOMA class
  \KOMAoptions{parskip=half}}
\makeatother
\usepackage{color}
\usepackage{fancyvrb}
\newcommand{\VerbBar}{|}
\newcommand{\VERB}{\Verb[commandchars=\\\{\}]}
\DefineVerbatimEnvironment{Highlighting}{Verbatim}{commandchars=\\\{\}}
% Add ',fontsize=\small' for more characters per line
\usepackage{framed}
\definecolor{shadecolor}{RGB}{248,248,248}
\newenvironment{Shaded}{\begin{snugshade}}{\end{snugshade}}
\newcommand{\AlertTok}[1]{\textcolor[rgb]{0.94,0.16,0.16}{#1}}
\newcommand{\AnnotationTok}[1]{\textcolor[rgb]{0.56,0.35,0.01}{\textbf{\textit{#1}}}}
\newcommand{\AttributeTok}[1]{\textcolor[rgb]{0.13,0.29,0.53}{#1}}
\newcommand{\BaseNTok}[1]{\textcolor[rgb]{0.00,0.00,0.81}{#1}}
\newcommand{\BuiltInTok}[1]{#1}
\newcommand{\CharTok}[1]{\textcolor[rgb]{0.31,0.60,0.02}{#1}}
\newcommand{\CommentTok}[1]{\textcolor[rgb]{0.56,0.35,0.01}{\textit{#1}}}
\newcommand{\CommentVarTok}[1]{\textcolor[rgb]{0.56,0.35,0.01}{\textbf{\textit{#1}}}}
\newcommand{\ConstantTok}[1]{\textcolor[rgb]{0.56,0.35,0.01}{#1}}
\newcommand{\ControlFlowTok}[1]{\textcolor[rgb]{0.13,0.29,0.53}{\textbf{#1}}}
\newcommand{\DataTypeTok}[1]{\textcolor[rgb]{0.13,0.29,0.53}{#1}}
\newcommand{\DecValTok}[1]{\textcolor[rgb]{0.00,0.00,0.81}{#1}}
\newcommand{\DocumentationTok}[1]{\textcolor[rgb]{0.56,0.35,0.01}{\textbf{\textit{#1}}}}
\newcommand{\ErrorTok}[1]{\textcolor[rgb]{0.64,0.00,0.00}{\textbf{#1}}}
\newcommand{\ExtensionTok}[1]{#1}
\newcommand{\FloatTok}[1]{\textcolor[rgb]{0.00,0.00,0.81}{#1}}
\newcommand{\FunctionTok}[1]{\textcolor[rgb]{0.13,0.29,0.53}{\textbf{#1}}}
\newcommand{\ImportTok}[1]{#1}
\newcommand{\InformationTok}[1]{\textcolor[rgb]{0.56,0.35,0.01}{\textbf{\textit{#1}}}}
\newcommand{\KeywordTok}[1]{\textcolor[rgb]{0.13,0.29,0.53}{\textbf{#1}}}
\newcommand{\NormalTok}[1]{#1}
\newcommand{\OperatorTok}[1]{\textcolor[rgb]{0.81,0.36,0.00}{\textbf{#1}}}
\newcommand{\OtherTok}[1]{\textcolor[rgb]{0.56,0.35,0.01}{#1}}
\newcommand{\PreprocessorTok}[1]{\textcolor[rgb]{0.56,0.35,0.01}{\textit{#1}}}
\newcommand{\RegionMarkerTok}[1]{#1}
\newcommand{\SpecialCharTok}[1]{\textcolor[rgb]{0.81,0.36,0.00}{\textbf{#1}}}
\newcommand{\SpecialStringTok}[1]{\textcolor[rgb]{0.31,0.60,0.02}{#1}}
\newcommand{\StringTok}[1]{\textcolor[rgb]{0.31,0.60,0.02}{#1}}
\newcommand{\VariableTok}[1]{\textcolor[rgb]{0.00,0.00,0.00}{#1}}
\newcommand{\VerbatimStringTok}[1]{\textcolor[rgb]{0.31,0.60,0.02}{#1}}
\newcommand{\WarningTok}[1]{\textcolor[rgb]{0.56,0.35,0.01}{\textbf{\textit{#1}}}}
\usepackage{graphicx}
\makeatletter
\newsavebox\pandoc@box
\newcommand*\pandocbounded[1]{% scales image to fit in text height/width
  \sbox\pandoc@box{#1}%
  \Gscale@div\@tempa{\textheight}{\dimexpr\ht\pandoc@box+\dp\pandoc@box\relax}%
  \Gscale@div\@tempb{\linewidth}{\wd\pandoc@box}%
  \ifdim\@tempb\p@<\@tempa\p@\let\@tempa\@tempb\fi% select the smaller of both
  \ifdim\@tempa\p@<\p@\scalebox{\@tempa}{\usebox\pandoc@box}%
  \else\usebox{\pandoc@box}%
  \fi%
}
% Set default figure placement to htbp
\def\fps@figure{htbp}
\makeatother
\setlength{\emergencystretch}{3em} % prevent overfull lines
\providecommand{\tightlist}{%
  \setlength{\itemsep}{0pt}\setlength{\parskip}{0pt}}
\usepackage{bookmark}
\IfFileExists{xurl.sty}{\usepackage{xurl}}{} % add URL line breaks if available
\urlstyle{same}
\hypersetup{
  pdftitle={Separation stats},
  pdfauthor={Sasha and Flavia},
  hidelinks,
  pdfcreator={LaTeX via pandoc}}

\title{Separation stats}
\author{Sasha and Flavia}
\date{}

\begin{document}
\maketitle

\section{Plot data}\label{plot-data}

\begin{Shaded}
\begin{Highlighting}[]
\FunctionTok{ggplot}\NormalTok{(isumm, }\FunctionTok{aes}\NormalTok{(treat, e.rel}\FloatTok{.130}\NormalTok{, }\AttributeTok{color =}\NormalTok{ trial)) }\SpecialCharTok{+} 
  \FunctionTok{geom\_point}\NormalTok{() }\SpecialCharTok{+} 
  \FunctionTok{geom\_line}\NormalTok{(}\FunctionTok{aes}\NormalTok{(}\AttributeTok{group =}\NormalTok{ trial)) }\SpecialCharTok{+} 
  \FunctionTok{facet\_wrap}\NormalTok{(}\SpecialCharTok{\textasciitilde{}}\NormalTok{ slurry.type, }\AttributeTok{scales =} \StringTok{\textquotesingle{}free\_x\textquotesingle{}}\NormalTok{) }\SpecialCharTok{+} 
  \FunctionTok{theme\_bw}\NormalTok{() }\SpecialCharTok{+} 
  \FunctionTok{labs}\NormalTok{(}\AttributeTok{y =} \StringTok{\textquotesingle{}Loss (frac. of TAN)\textquotesingle{}}\NormalTok{) }\SpecialCharTok{+} 
  \FunctionTok{theme}\NormalTok{(}\AttributeTok{legend.title =} \FunctionTok{element\_blank}\NormalTok{())}
\end{Highlighting}
\end{Shaded}

\pandocbounded{\includegraphics[keepaspectratio]{stats_files/figure-latex/unnamed-chunk-1-1.pdf}}

Through means.

\begin{Shaded}
\begin{Highlighting}[]
\FunctionTok{ggplot}\NormalTok{(isumm, }\FunctionTok{aes}\NormalTok{(treat, e.rel}\FloatTok{.130}\NormalTok{, }\AttributeTok{color =}\NormalTok{ trial)) }\SpecialCharTok{+} 
  \FunctionTok{geom\_line}\NormalTok{() }\SpecialCharTok{+} 
  \FunctionTok{geom\_point}\NormalTok{() }\SpecialCharTok{+} 
  \FunctionTok{stat\_summary}\NormalTok{(}\FunctionTok{aes}\NormalTok{(}\AttributeTok{group =}\NormalTok{ trial), }\AttributeTok{fun =}\NormalTok{ mean, }\AttributeTok{geom =} \StringTok{\textquotesingle{}line\textquotesingle{}}\NormalTok{) }\SpecialCharTok{+}
  \FunctionTok{facet\_wrap}\NormalTok{(}\SpecialCharTok{\textasciitilde{}}\NormalTok{ slurry.type, }\AttributeTok{scales =} \StringTok{\textquotesingle{}free\_x\textquotesingle{}}\NormalTok{) }\SpecialCharTok{+} 
  \FunctionTok{theme\_bw}\NormalTok{() }\SpecialCharTok{+} 
  \FunctionTok{labs}\NormalTok{(}\AttributeTok{y =} \StringTok{\textquotesingle{}Loss (frac. of TAN)\textquotesingle{}}\NormalTok{) }\SpecialCharTok{+} 
  \FunctionTok{theme}\NormalTok{(}\AttributeTok{legend.title =} \FunctionTok{element\_blank}\NormalTok{())}
\end{Highlighting}
\end{Shaded}

\pandocbounded{\includegraphics[keepaspectratio]{stats_files/figure-latex/unnamed-chunk-2-1.pdf}}

\begin{Shaded}
\begin{Highlighting}[]
\FunctionTok{ggplot}\NormalTok{(isumm, }\FunctionTok{aes}\NormalTok{(}\FunctionTok{as.integer}\NormalTok{(}\FunctionTok{factor}\NormalTok{(treat)), e.rel}\FloatTok{.130}\NormalTok{, }\AttributeTok{color =}\NormalTok{ trial)) }\SpecialCharTok{+} 
  \FunctionTok{geom\_point}\NormalTok{() }\SpecialCharTok{+} 
  \FunctionTok{geom\_smooth}\NormalTok{(}\AttributeTok{se =} \ConstantTok{FALSE}\NormalTok{) }\SpecialCharTok{+}
  \FunctionTok{facet\_wrap}\NormalTok{(}\SpecialCharTok{\textasciitilde{}}\NormalTok{ slurry.type, }\AttributeTok{scales =} \StringTok{\textquotesingle{}free\_x\textquotesingle{}}\NormalTok{) }\SpecialCharTok{+} 
  \FunctionTok{theme\_bw}\NormalTok{() }\SpecialCharTok{+} 
  \FunctionTok{labs}\NormalTok{(}\AttributeTok{y =} \StringTok{\textquotesingle{}Loss (frac. of TAN)\textquotesingle{}}\NormalTok{) }\SpecialCharTok{+} 
  \FunctionTok{theme}\NormalTok{(}\AttributeTok{legend.title =} \FunctionTok{element\_blank}\NormalTok{())}
\end{Highlighting}
\end{Shaded}

\begin{verbatim}
## `geom_smooth()` using method = 'loess' and formula = 'y ~ x'
\end{verbatim}

\begin{verbatim}
## Warning in simpleLoess(y, x, w, span, degree = degree, parametric =
## parametric, : pseudoinverse used at 0.99
\end{verbatim}

\begin{verbatim}
## Warning in simpleLoess(y, x, w, span, degree = degree, parametric =
## parametric, : neighborhood radius 1.01
\end{verbatim}

\begin{verbatim}
## Warning in simpleLoess(y, x, w, span, degree = degree, parametric =
## parametric, : reciprocal condition number 0
\end{verbatim}

\begin{verbatim}
## Warning in simpleLoess(y, x, w, span, degree = degree, parametric =
## parametric, : There are other near singularities as well. 1.0201
\end{verbatim}

\begin{verbatim}
## Warning in simpleLoess(y, x, w, span, degree = degree, parametric =
## parametric, : pseudoinverse used at 0.99
\end{verbatim}

\begin{verbatim}
## Warning in simpleLoess(y, x, w, span, degree = degree, parametric =
## parametric, : neighborhood radius 1.01
\end{verbatim}

\begin{verbatim}
## Warning in simpleLoess(y, x, w, span, degree = degree, parametric =
## parametric, : reciprocal condition number 0
\end{verbatim}

\begin{verbatim}
## Warning in simpleLoess(y, x, w, span, degree = degree, parametric =
## parametric, : There are other near singularities as well. 1.0201
\end{verbatim}

\begin{verbatim}
## Warning in simpleLoess(y, x, w, span, degree = degree, parametric =
## parametric, : pseudoinverse used at 0.99
\end{verbatim}

\begin{verbatim}
## Warning in simpleLoess(y, x, w, span, degree = degree, parametric =
## parametric, : neighborhood radius 2.01
\end{verbatim}

\begin{verbatim}
## Warning in simpleLoess(y, x, w, span, degree = degree, parametric =
## parametric, : reciprocal condition number 0
\end{verbatim}

\begin{verbatim}
## Warning in simpleLoess(y, x, w, span, degree = degree, parametric =
## parametric, : There are other near singularities as well. 1.0201
\end{verbatim}

\begin{verbatim}
## Warning in simpleLoess(y, x, w, span, degree = degree, parametric =
## parametric, : pseudoinverse used at 3.99
\end{verbatim}

\begin{verbatim}
## Warning in simpleLoess(y, x, w, span, degree = degree, parametric =
## parametric, : neighborhood radius 1.01
\end{verbatim}

\begin{verbatim}
## Warning in simpleLoess(y, x, w, span, degree = degree, parametric =
## parametric, : reciprocal condition number 0
\end{verbatim}

\begin{verbatim}
## Warning in simpleLoess(y, x, w, span, degree = degree, parametric =
## parametric, : There are other near singularities as well. 1.0201
\end{verbatim}

\begin{verbatim}
## Warning in simpleLoess(y, x, w, span, degree = degree, parametric =
## parametric, : pseudoinverse used at 3.99
\end{verbatim}

\begin{verbatim}
## Warning in simpleLoess(y, x, w, span, degree = degree, parametric =
## parametric, : neighborhood radius 1.01
\end{verbatim}

\begin{verbatim}
## Warning in simpleLoess(y, x, w, span, degree = degree, parametric =
## parametric, : reciprocal condition number 0
\end{verbatim}

\begin{verbatim}
## Warning in simpleLoess(y, x, w, span, degree = degree, parametric =
## parametric, : There are other near singularities as well. 1.0201
\end{verbatim}

\begin{verbatim}
## Warning in simpleLoess(y, x, w, span, degree = degree, parametric =
## parametric, : pseudoinverse used at 3.99
\end{verbatim}

\begin{verbatim}
## Warning in simpleLoess(y, x, w, span, degree = degree, parametric =
## parametric, : neighborhood radius 1.01
\end{verbatim}

\begin{verbatim}
## Warning in simpleLoess(y, x, w, span, degree = degree, parametric =
## parametric, : reciprocal condition number 0
\end{verbatim}

\begin{verbatim}
## Warning in simpleLoess(y, x, w, span, degree = degree, parametric =
## parametric, : There are other near singularities as well. 1.0201
\end{verbatim}

\pandocbounded{\includegraphics[keepaspectratio]{stats_files/figure-latex/unnamed-chunk-3-1.pdf}}

\begin{Shaded}
\begin{Highlighting}[]
\FunctionTok{ggplot}\NormalTok{(isumm, }\FunctionTok{aes}\NormalTok{(treat, e.rel}\FloatTok{.130}\NormalTok{, }\AttributeTok{color =}\NormalTok{ trial)) }\SpecialCharTok{+} 
  \FunctionTok{geom\_point}\NormalTok{() }\SpecialCharTok{+} 
  \FunctionTok{geom\_line}\NormalTok{() }\SpecialCharTok{+} 
  \FunctionTok{geom\_smooth}\NormalTok{(}\AttributeTok{se =} \ConstantTok{FALSE}\NormalTok{) }\SpecialCharTok{+}
  \FunctionTok{facet\_wrap}\NormalTok{(}\SpecialCharTok{\textasciitilde{}}\NormalTok{ slurry.type, }\AttributeTok{scales =} \StringTok{\textquotesingle{}free\_x\textquotesingle{}}\NormalTok{) }\SpecialCharTok{+} 
  \FunctionTok{theme\_bw}\NormalTok{() }\SpecialCharTok{+} 
  \FunctionTok{labs}\NormalTok{(}\AttributeTok{y =} \StringTok{\textquotesingle{}Loss (frac. of TAN)\textquotesingle{}}\NormalTok{) }\SpecialCharTok{+} 
  \FunctionTok{theme}\NormalTok{(}\AttributeTok{legend.title =} \FunctionTok{element\_blank}\NormalTok{())}
\end{Highlighting}
\end{Shaded}

\begin{verbatim}
## `geom_smooth()` using method = 'loess' and formula = 'y ~ x'
\end{verbatim}

\pandocbounded{\includegraphics[keepaspectratio]{stats_files/figure-latex/unnamed-chunk-4-1.pdf}}

\section{Mixed-effects models}\label{mixed-effects-models}

\subsection{Model 1}\label{model-1}

First ignore any variation in separation response among trials (plot
shows evidence that we should \emph{not} ignore it).

\begin{Shaded}
\begin{Highlighting}[]
\NormalTok{mod1 }\OtherTok{\textless{}{-}} \FunctionTok{lmer}\NormalTok{(e.rel}\FloatTok{.130} \SpecialCharTok{\textasciitilde{}}\NormalTok{ separation }\SpecialCharTok{*}\NormalTok{ slurry.type }\SpecialCharTok{+}\NormalTok{ (}\DecValTok{1}\SpecialCharTok{|}\NormalTok{trial), }\AttributeTok{data =}\NormalTok{ isumm, }\AttributeTok{REML =} \ConstantTok{FALSE}\NormalTok{)}
\end{Highlighting}
\end{Shaded}

\begin{verbatim}
## boundary (singular) fit: see help('isSingular')
\end{verbatim}

\begin{Shaded}
\begin{Highlighting}[]
\FunctionTok{summary}\NormalTok{(mod1)}
\end{Highlighting}
\end{Shaded}

\begin{verbatim}
## Linear mixed model fit by maximum likelihood  ['lmerMod']
## Formula: e.rel.130 ~ separation * slurry.type + (1 | trial)
##    Data: isumm
## 
##       AIC       BIC    logLik -2*log(L)  df.resid 
##     -99.3     -83.5      57.6    -115.3        45 
## 
## Scaled residuals: 
##      Min       1Q   Median       3Q      Max 
## -2.10599 -0.59951 -0.06593  0.64664  2.46112 
## 
## Random effects:
##  Groups   Name        Variance  Std.Dev. 
##  trial    (Intercept) 1.145e-21 3.384e-11
##  Residual             6.653e-03 8.157e-02
## Number of obs: 53, groups:  trial, 6
## 
## Fixed effects:
##                                Estimate Std. Error t value
## (Intercept)                     0.42400    0.02884  14.703
## separationLiquid               -0.24602    0.03963  -6.207
## separationSolid                -0.36210    0.03963  -9.136
## slurry.typePS                  -0.19174    0.03963  -4.838
## separationLiquid:slurry.typePS  0.06774    0.05522   1.227
## separationSolid:slurry.typePS   0.49762    0.05522   9.012
## 
## Correlation of Fixed Effects:
##             (Intr) sprtnL sprtnS slr.PS sL:.PS
## separatnLqd -0.728                            
## separatnSld -0.728  0.529                     
## slrry.typPS -0.728  0.529  0.529              
## sprtnLq:.PS  0.522 -0.718 -0.380 -0.718       
## sprtnSl:.PS  0.522 -0.380 -0.718 -0.718  0.515
## optimizer (nloptwrap) convergence code: 0 (OK)
## boundary (singular) fit: see help('isSingular')
\end{verbatim}

\begin{Shaded}
\begin{Highlighting}[]
\FunctionTok{anova}\NormalTok{(mod1)}
\end{Highlighting}
\end{Shaded}

\begin{verbatim}
## Analysis of Variance Table
##                        npar  Sum Sq Mean Sq F value
## separation                2 0.37288 0.18644 28.0230
## slurry.type               1 0.00000 0.00000  0.0005
## separation:slurry.type    2 0.64593 0.32296 48.5437
\end{verbatim}

\begin{Shaded}
\begin{Highlighting}[]
\FunctionTok{plot}\NormalTok{(mod1)}
\end{Highlighting}
\end{Shaded}

\pandocbounded{\includegraphics[keepaspectratio]{stats_files/figure-latex/unnamed-chunk-6-1.pdf}}

\begin{Shaded}
\begin{Highlighting}[]
\FunctionTok{qqnorm}\NormalTok{(}\FunctionTok{resid}\NormalTok{(mod1))}
\FunctionTok{qqline}\NormalTok{(}\FunctionTok{resid}\NormalTok{(mod1))}
\end{Highlighting}
\end{Shaded}

\pandocbounded{\includegraphics[keepaspectratio]{stats_files/figure-latex/unnamed-chunk-6-2.pdf}}

\subsection{Model 2}\label{model-2}

Now include this random separation effect across trials.

\begin{Shaded}
\begin{Highlighting}[]
\NormalTok{mod2 }\OtherTok{\textless{}{-}} \FunctionTok{lmer}\NormalTok{(e.rel}\FloatTok{.130} \SpecialCharTok{\textasciitilde{}}\NormalTok{ separation }\SpecialCharTok{*}\NormalTok{ slurry.type }\SpecialCharTok{+}\NormalTok{ (}\DecValTok{1} \SpecialCharTok{+}\NormalTok{ separation }\SpecialCharTok{|}\NormalTok{ trial), }\AttributeTok{data =}\NormalTok{ isumm, }\AttributeTok{REML =} \ConstantTok{FALSE}\NormalTok{)}
\FunctionTok{summary}\NormalTok{(mod2)}
\end{Highlighting}
\end{Shaded}

\begin{verbatim}
## Linear mixed model fit by maximum likelihood  ['lmerMod']
## Formula: e.rel.130 ~ separation * slurry.type + (1 + separation | trial)
##    Data: isumm
## 
##       AIC       BIC    logLik -2*log(L)  df.resid 
##    -128.2    -102.6      77.1    -154.2        40 
## 
## Scaled residuals: 
##      Min       1Q   Median       3Q      Max 
## -1.61123 -0.52248 -0.04807  0.57283  2.47528 
## 
## Random effects:
##  Groups   Name             Variance Std.Dev. Corr       
##  trial    (Intercept)      0.003554 0.05961             
##           separationLiquid 0.008350 0.09138  -0.88      
##           separationSolid  0.013958 0.11815  -0.63  0.22
##  Residual                  0.001752 0.04185             
## Number of obs: 53, groups:  trial, 6
## 
## Fixed effects:
##                                Estimate Std. Error t value
## (Intercept)                     0.42071    0.03754  11.207
## separationLiquid               -0.24273    0.05659  -4.289
## separationSolid                -0.35881    0.07122  -5.038
## slurry.typePS                  -0.18846    0.05281  -3.569
## separationLiquid:slurry.typePS  0.06446    0.07985   0.807
## separationSolid:slurry.typePS   0.49433    0.10057   4.915
## 
## Correlation of Fixed Effects:
##             (Intr) sprtnL sprtnS slr.PS sL:.PS
## separatnLqd -0.859                            
## separatnSld -0.636  0.252                     
## slrry.typPS -0.711  0.611  0.452              
## sprtnLq:.PS  0.609 -0.709 -0.179 -0.858       
## sprtnSl:.PS  0.451 -0.179 -0.708 -0.635  0.250
\end{verbatim}

\begin{Shaded}
\begin{Highlighting}[]
\FunctionTok{anova}\NormalTok{(mod2)}
\end{Highlighting}
\end{Shaded}

\begin{verbatim}
## Analysis of Variance Table
##                        npar   Sum Sq  Mean Sq F value
## separation                2 0.049490 0.024745 14.1263
## slurry.type               1 0.013629 0.013629  7.7807
## separation:slurry.type    2 0.042650 0.021325 12.1740
\end{verbatim}

\subsection{Model 2 hypothsis tests}\label{model-2-hypothsis-tests}

We can determine if there is evidence for it with a likelihood ratio
test.

\begin{Shaded}
\begin{Highlighting}[]
\FunctionTok{anova}\NormalTok{(mod1, mod2, }\AttributeTok{test =} \StringTok{\textquotesingle{}Chisq\textquotesingle{}}\NormalTok{)}
\end{Highlighting}
\end{Shaded}

\begin{verbatim}
## Data: isumm
## Models:
## mod1: e.rel.130 ~ separation * slurry.type + (1 | trial)
## mod2: e.rel.130 ~ separation * slurry.type + (1 + separation | trial)
##      npar      AIC      BIC logLik -2*log(L)  Chisq Df Pr(>Chisq)    
## mod1    8  -99.264  -83.502 57.632   -115.26                         
## mod2   13 -128.247 -102.633 77.124   -154.25 38.983  5  2.394e-07 ***
## ---
## Signif. codes:  0 '***' 0.001 '**' 0.01 '*' 0.05 '.' 0.1 ' ' 1
\end{verbatim}

Strong evidence, with p value below 1E-6.

For interaction term:

\begin{Shaded}
\begin{Highlighting}[]
\NormalTok{mod2b }\OtherTok{\textless{}{-}} \FunctionTok{lmer}\NormalTok{(e.rel}\FloatTok{.130} \SpecialCharTok{\textasciitilde{}}\NormalTok{ separation }\SpecialCharTok{+}\NormalTok{ slurry.type }\SpecialCharTok{+}\NormalTok{ (}\DecValTok{1} \SpecialCharTok{+}\NormalTok{ separation }\SpecialCharTok{|}\NormalTok{ trial), }\AttributeTok{data =}\NormalTok{ isumm, }\AttributeTok{REML =} \ConstantTok{FALSE}\NormalTok{)}
\FunctionTok{anova}\NormalTok{(mod2, mod2b, }\AttributeTok{test =} \StringTok{\textquotesingle{}Chisq\textquotesingle{}}\NormalTok{)}
\end{Highlighting}
\end{Shaded}

\begin{verbatim}
## Data: isumm
## Models:
## mod2b: e.rel.130 ~ separation + slurry.type + (1 + separation | trial)
## mod2: e.rel.130 ~ separation * slurry.type + (1 + separation | trial)
##       npar     AIC     BIC logLik -2*log(L)  Chisq Df Pr(>Chisq)   
## mod2b   11 -122.51 -100.84 72.255   -144.51                        
## mod2    13 -128.25 -102.63 77.124   -154.25 9.7378  2   0.007682 **
## ---
## Signif. codes:  0 '***' 0.001 '**' 0.01 '*' 0.05 '.' 0.1 ' ' 1
\end{verbatim}

p \textless{} 0.01, pretty clear.

And it is clear from the plots that there is no consistent difference
between

\begin{itemize}
\tightlist
\item
  solid and liquid fractions from AD,
\item
  raw and solid from pig slurry.
\end{itemize}

The only thing not clear from the combination of these results and the
plots is if.

\begin{itemize}
\tightlist
\item
  AD liquid fraction is in fact lower than raw AD,
\item
  pig liquid fraction is in fact lower than raw pig slurry.
\end{itemize}

We can test exactly these and only these by grouping factor levels and
using likelihood ratio tests.

\begin{Shaded}
\begin{Highlighting}[]
\NormalTok{isumm[, separation.psg }\SpecialCharTok{:=}\NormalTok{ separation]}
\NormalTok{isumm[slurry.type }\SpecialCharTok{==} \StringTok{\textquotesingle{}PS\textquotesingle{}} \SpecialCharTok{\&}\NormalTok{ separation }\SpecialCharTok{==} \StringTok{\textquotesingle{}Liquid\textquotesingle{}}\NormalTok{, separation.psg }\SpecialCharTok{:=} \StringTok{\textquotesingle{}Unsep\textquotesingle{}}\NormalTok{]}
\FunctionTok{table}\NormalTok{(isumm[, .(slurry.type, separation, trial)])}
\end{Highlighting}
\end{Shaded}

\begin{verbatim}
## , , trial = Trial 1
## 
##            separation
## slurry.type Unsep Liquid Solid
##          AD     3      3     3
##          PS     0      0     0
## 
## , , trial = Trial 2
## 
##            separation
## slurry.type Unsep Liquid Solid
##          AD     0      0     0
##          PS     3      3     3
## 
## , , trial = Trial 3
## 
##            separation
## slurry.type Unsep Liquid Solid
##          AD     3      3     3
##          PS     0      0     0
## 
## , , trial = Trial 4
## 
##            separation
## slurry.type Unsep Liquid Solid
##          AD     0      0     0
##          PS     3      3     3
## 
## , , trial = Trial 5
## 
##            separation
## slurry.type Unsep Liquid Solid
##          AD     2      3     3
##          PS     0      0     0
## 
## , , trial = Trial 6
## 
##            separation
## slurry.type Unsep Liquid Solid
##          AD     0      0     0
##          PS     3      3     3
\end{verbatim}

\begin{Shaded}
\begin{Highlighting}[]
\FunctionTok{table}\NormalTok{(isumm[, .(slurry.type, separation.psg, trial)])}
\end{Highlighting}
\end{Shaded}

\begin{verbatim}
## , , trial = Trial 1
## 
##            separation.psg
## slurry.type Unsep Liquid Solid
##          AD     3      3     3
##          PS     0      0     0
## 
## , , trial = Trial 2
## 
##            separation.psg
## slurry.type Unsep Liquid Solid
##          AD     0      0     0
##          PS     6      0     3
## 
## , , trial = Trial 3
## 
##            separation.psg
## slurry.type Unsep Liquid Solid
##          AD     3      3     3
##          PS     0      0     0
## 
## , , trial = Trial 4
## 
##            separation.psg
## slurry.type Unsep Liquid Solid
##          AD     0      0     0
##          PS     6      0     3
## 
## , , trial = Trial 5
## 
##            separation.psg
## slurry.type Unsep Liquid Solid
##          AD     2      3     3
##          PS     0      0     0
## 
## , , trial = Trial 6
## 
##            separation.psg
## slurry.type Unsep Liquid Solid
##          AD     0      0     0
##          PS     6      0     3
\end{verbatim}

\begin{Shaded}
\begin{Highlighting}[]
\NormalTok{isumm[, separation.adg }\SpecialCharTok{:=}\NormalTok{ separation]}
\NormalTok{isumm[slurry.type }\SpecialCharTok{==} \StringTok{\textquotesingle{}AD\textquotesingle{}} \SpecialCharTok{\&}\NormalTok{ separation }\SpecialCharTok{==} \StringTok{\textquotesingle{}Liquid\textquotesingle{}}\NormalTok{, separation.psg }\SpecialCharTok{:=} \StringTok{\textquotesingle{}Unsep\textquotesingle{}}\NormalTok{]}
\FunctionTok{table}\NormalTok{(isumm[, .(slurry.type, separation, trial)])}
\end{Highlighting}
\end{Shaded}

\begin{verbatim}
## , , trial = Trial 1
## 
##            separation
## slurry.type Unsep Liquid Solid
##          AD     3      3     3
##          PS     0      0     0
## 
## , , trial = Trial 2
## 
##            separation
## slurry.type Unsep Liquid Solid
##          AD     0      0     0
##          PS     3      3     3
## 
## , , trial = Trial 3
## 
##            separation
## slurry.type Unsep Liquid Solid
##          AD     3      3     3
##          PS     0      0     0
## 
## , , trial = Trial 4
## 
##            separation
## slurry.type Unsep Liquid Solid
##          AD     0      0     0
##          PS     3      3     3
## 
## , , trial = Trial 5
## 
##            separation
## slurry.type Unsep Liquid Solid
##          AD     2      3     3
##          PS     0      0     0
## 
## , , trial = Trial 6
## 
##            separation
## slurry.type Unsep Liquid Solid
##          AD     0      0     0
##          PS     3      3     3
\end{verbatim}

\begin{Shaded}
\begin{Highlighting}[]
\FunctionTok{table}\NormalTok{(isumm[, .(slurry.type, separation.adg, trial)])}
\end{Highlighting}
\end{Shaded}

\begin{verbatim}
## , , trial = Trial 1
## 
##            separation.adg
## slurry.type Unsep Liquid Solid
##          AD     3      3     3
##          PS     0      0     0
## 
## , , trial = Trial 2
## 
##            separation.adg
## slurry.type Unsep Liquid Solid
##          AD     0      0     0
##          PS     3      3     3
## 
## , , trial = Trial 3
## 
##            separation.adg
## slurry.type Unsep Liquid Solid
##          AD     3      3     3
##          PS     0      0     0
## 
## , , trial = Trial 4
## 
##            separation.adg
## slurry.type Unsep Liquid Solid
##          AD     0      0     0
##          PS     3      3     3
## 
## , , trial = Trial 5
## 
##            separation.adg
## slurry.type Unsep Liquid Solid
##          AD     2      3     3
##          PS     0      0     0
## 
## , , trial = Trial 6
## 
##            separation.adg
## slurry.type Unsep Liquid Solid
##          AD     0      0     0
##          PS     3      3     3
\end{verbatim}

\begin{Shaded}
\begin{Highlighting}[]
\NormalTok{mod2c }\OtherTok{\textless{}{-}} \FunctionTok{lmer}\NormalTok{(e.rel}\FloatTok{.130} \SpecialCharTok{\textasciitilde{}}\NormalTok{ separation.psg }\SpecialCharTok{+}\NormalTok{ slurry.type }\SpecialCharTok{+}\NormalTok{ (}\DecValTok{1} \SpecialCharTok{+}\NormalTok{ separation.psg }\SpecialCharTok{|}\NormalTok{ trial), }\AttributeTok{data =}\NormalTok{ isumm, }\AttributeTok{REML =} \ConstantTok{FALSE}\NormalTok{)}
\end{Highlighting}
\end{Shaded}

\begin{verbatim}
## boundary (singular) fit: see help('isSingular')
\end{verbatim}

\begin{Shaded}
\begin{Highlighting}[]
\NormalTok{mod2d }\OtherTok{\textless{}{-}} \FunctionTok{lmer}\NormalTok{(e.rel}\FloatTok{.130} \SpecialCharTok{\textasciitilde{}}\NormalTok{ separation.adg }\SpecialCharTok{+}\NormalTok{ slurry.type }\SpecialCharTok{+}\NormalTok{ (}\DecValTok{1} \SpecialCharTok{+}\NormalTok{ separation.adg }\SpecialCharTok{|}\NormalTok{ trial), }\AttributeTok{data =}\NormalTok{ isumm, }\AttributeTok{REML =} \ConstantTok{FALSE}\NormalTok{)}
\FunctionTok{anova}\NormalTok{(mod2, mod2c, }\AttributeTok{test =} \StringTok{\textquotesingle{}Chisq\textquotesingle{}}\NormalTok{)}
\end{Highlighting}
\end{Shaded}

\begin{verbatim}
## Data: isumm
## Models:
## mod2c: e.rel.130 ~ separation.psg + slurry.type + (1 + separation.psg | trial)
## mod2: e.rel.130 ~ separation * slurry.type + (1 + separation | trial)
##       npar      AIC      BIC logLik -2*log(L)  Chisq Df Pr(>Chisq)    
## mod2c    7  -54.794  -41.002 34.397   -68.794                         
## mod2    13 -128.247 -102.633 77.124  -154.247 85.454  6  2.659e-16 ***
## ---
## Signif. codes:  0 '***' 0.001 '**' 0.01 '*' 0.05 '.' 0.1 ' ' 1
\end{verbatim}

\begin{Shaded}
\begin{Highlighting}[]
\FunctionTok{anova}\NormalTok{(mod2, mod2d, }\AttributeTok{test =} \StringTok{\textquotesingle{}Chisq\textquotesingle{}}\NormalTok{)}
\end{Highlighting}
\end{Shaded}

\begin{verbatim}
## Data: isumm
## Models:
## mod2d: e.rel.130 ~ separation.adg + slurry.type + (1 + separation.adg | trial)
## mod2: e.rel.130 ~ separation * slurry.type + (1 + separation | trial)
##       npar     AIC     BIC logLik -2*log(L)  Chisq Df Pr(>Chisq)   
## mod2d   11 -122.51 -100.84 72.255   -144.51                        
## mod2    13 -128.25 -102.63 77.124   -154.25 9.7378  2   0.007682 **
## ---
## Signif. codes:  0 '***' 0.001 '**' 0.01 '*' 0.05 '.' 0.1 ' ' 1
\end{verbatim}

Singular fit is problematic for comparison.

\begin{Shaded}
\begin{Highlighting}[]
\NormalTok{mod1c }\OtherTok{\textless{}{-}} \FunctionTok{lmer}\NormalTok{(e.rel}\FloatTok{.130} \SpecialCharTok{\textasciitilde{}}\NormalTok{ separation.psg }\SpecialCharTok{+}\NormalTok{ slurry.type }\SpecialCharTok{+}\NormalTok{ (}\DecValTok{1} \SpecialCharTok{|}\NormalTok{ trial), }\AttributeTok{data =}\NormalTok{ isumm, }\AttributeTok{REML =} \ConstantTok{FALSE}\NormalTok{)}
\end{Highlighting}
\end{Shaded}

\begin{verbatim}
## boundary (singular) fit: see help('isSingular')
\end{verbatim}

\begin{Shaded}
\begin{Highlighting}[]
\NormalTok{mod1d }\OtherTok{\textless{}{-}} \FunctionTok{lmer}\NormalTok{(e.rel}\FloatTok{.130} \SpecialCharTok{\textasciitilde{}}\NormalTok{ separation.adg }\SpecialCharTok{+}\NormalTok{ slurry.type }\SpecialCharTok{+}\NormalTok{ (}\DecValTok{1} \SpecialCharTok{|}\NormalTok{ trial), }\AttributeTok{data =}\NormalTok{ isumm, }\AttributeTok{REML =} \ConstantTok{FALSE}\NormalTok{)}
\end{Highlighting}
\end{Shaded}

\begin{verbatim}
## boundary (singular) fit: see help('isSingular')
\end{verbatim}

\begin{Shaded}
\begin{Highlighting}[]
\FunctionTok{anova}\NormalTok{(mod1, mod1c, }\AttributeTok{test =} \StringTok{\textquotesingle{}Chisq\textquotesingle{}}\NormalTok{)}
\end{Highlighting}
\end{Shaded}

\begin{verbatim}
## Data: isumm
## Models:
## mod1c: e.rel.130 ~ separation.psg + slurry.type + (1 | trial)
## mod1: e.rel.130 ~ separation * slurry.type + (1 | trial)
##       npar     AIC     BIC logLik -2*log(L)  Chisq Df Pr(>Chisq)    
## mod1c    5 -33.289 -23.438 21.645   -43.289                         
## mod1     8 -99.264 -83.502 57.632  -115.264 71.975  3  1.612e-15 ***
## ---
## Signif. codes:  0 '***' 0.001 '**' 0.01 '*' 0.05 '.' 0.1 ' ' 1
\end{verbatim}

\begin{Shaded}
\begin{Highlighting}[]
\FunctionTok{anova}\NormalTok{(mod1, mod1d, }\AttributeTok{test =} \StringTok{\textquotesingle{}Chisq\textquotesingle{}}\NormalTok{)}
\end{Highlighting}
\end{Shaded}

\begin{verbatim}
## Data: isumm
## Models:
## mod1d: e.rel.130 ~ separation.adg + slurry.type + (1 | trial)
## mod1: e.rel.130 ~ separation * slurry.type + (1 | trial)
##       npar     AIC     BIC logLik -2*log(L)  Chisq Df Pr(>Chisq)    
## mod1d    6 -48.095 -36.274 30.048   -60.095                         
## mod1     8 -99.264 -83.502 57.632  -115.264 55.169  2  1.048e-12 ***
## ---
## Signif. codes:  0 '***' 0.001 '**' 0.01 '*' 0.05 '.' 0.1 ' ' 1
\end{verbatim}

No better. Try subset instead.

\begin{Shaded}
\begin{Highlighting}[]
\NormalTok{isummul }\OtherTok{\textless{}{-}}\NormalTok{ isumm[separation }\SpecialCharTok{\%in\%} \FunctionTok{c}\NormalTok{(}\StringTok{\textquotesingle{}Unsep\textquotesingle{}}\NormalTok{, }\StringTok{\textquotesingle{}Liquid\textquotesingle{}}\NormalTok{)]}
\end{Highlighting}
\end{Shaded}

\begin{Shaded}
\begin{Highlighting}[]
\NormalTok{mod3 }\OtherTok{\textless{}{-}} \FunctionTok{lmer}\NormalTok{(e.rel}\FloatTok{.130} \SpecialCharTok{\textasciitilde{}}\NormalTok{ separation }\SpecialCharTok{*}\NormalTok{ slurry.type }\SpecialCharTok{+}\NormalTok{ (}\DecValTok{1}\SpecialCharTok{|}\NormalTok{trial), }\AttributeTok{data =}\NormalTok{ isummul, }\AttributeTok{REML =} \ConstantTok{FALSE}\NormalTok{)}
\FunctionTok{summary}\NormalTok{(mod3)}
\end{Highlighting}
\end{Shaded}

\begin{verbatim}
## Linear mixed model fit by maximum likelihood  ['lmerMod']
## Formula: e.rel.130 ~ separation * slurry.type + (1 | trial)
##    Data: isummul
## 
##       AIC       BIC    logLik -2*log(L)  df.resid 
##     -76.2     -66.9      44.1     -88.2        29 
## 
## Scaled residuals: 
##      Min       1Q   Median       3Q      Max 
## -1.93249 -0.72894 -0.08274  0.56242  2.82139 
## 
## Random effects:
##  Groups   Name        Variance  Std.Dev.
##  trial    (Intercept) 0.0005319 0.02306 
##  Residual             0.0042883 0.06548 
## Number of obs: 35, groups:  trial, 6
## 
## Fixed effects:
##                                Estimate Std. Error t value
## (Intercept)                     0.42616    0.02677  15.919
## separationLiquid               -0.24818    0.03187  -7.787
## slurry.typePS                  -0.19390    0.03702  -5.238
## separationLiquid:slurry.typePS  0.06990    0.04437   1.575
## 
## Correlation of Fixed Effects:
##             (Intr) sprtnL slr.PS
## separatnLqd -0.632              
## slrry.typPS -0.723  0.457       
## sprtnLq:.PS  0.454 -0.718 -0.618
\end{verbatim}

\begin{Shaded}
\begin{Highlighting}[]
\FunctionTok{anova}\NormalTok{(mod3)}
\end{Highlighting}
\end{Shaded}

\begin{verbatim}
## Analysis of Variance Table
##                        npar  Sum Sq Mean Sq F value
## separation                1 0.38217 0.38217  89.119
## slurry.type               1 0.12623 0.12623  29.436
## separation:slurry.type    1 0.01064 0.01064   2.482
\end{verbatim}

Check for interaction now.

\begin{Shaded}
\begin{Highlighting}[]
\NormalTok{mod3b }\OtherTok{\textless{}{-}} \FunctionTok{lmer}\NormalTok{(e.rel}\FloatTok{.130} \SpecialCharTok{\textasciitilde{}}\NormalTok{ separation }\SpecialCharTok{+}\NormalTok{ slurry.type }\SpecialCharTok{+}\NormalTok{ (}\DecValTok{1}\SpecialCharTok{|}\NormalTok{trial), }\AttributeTok{data =}\NormalTok{ isummul, }\AttributeTok{REML =} \ConstantTok{FALSE}\NormalTok{)}
\FunctionTok{anova}\NormalTok{(mod3, mod3b, }\AttributeTok{test =} \StringTok{\textquotesingle{}Chisq\textquotesingle{}}\NormalTok{)}
\end{Highlighting}
\end{Shaded}

\begin{verbatim}
## Data: isummul
## Models:
## mod3b: e.rel.130 ~ separation + slurry.type + (1 | trial)
## mod3: e.rel.130 ~ separation * slurry.type + (1 | trial)
##       npar     AIC     BIC logLik -2*log(L)  Chisq Df Pr(>Chisq)
## mod3b    5 -75.866 -68.090 42.933   -85.866                     
## mod3     6 -76.226 -66.894 44.113   -88.226 2.3593  1     0.1245
\end{verbatim}

Nope! So solid fraction response drives interaction. That is obvious
from the plots.

So we can easily test for overall unsep vs.~liquid effect.

\begin{Shaded}
\begin{Highlighting}[]
\NormalTok{mod3c }\OtherTok{\textless{}{-}} \FunctionTok{lmer}\NormalTok{(e.rel}\FloatTok{.130} \SpecialCharTok{\textasciitilde{}}\NormalTok{ slurry.type }\SpecialCharTok{+}\NormalTok{ (}\DecValTok{1}\SpecialCharTok{|}\NormalTok{trial), }\AttributeTok{data =}\NormalTok{ isummul, }\AttributeTok{REML =} \ConstantTok{FALSE}\NormalTok{)}
\end{Highlighting}
\end{Shaded}

\begin{verbatim}
## boundary (singular) fit: see help('isSingular')
\end{verbatim}

\begin{Shaded}
\begin{Highlighting}[]
\FunctionTok{summary}\NormalTok{(mod3c)}
\end{Highlighting}
\end{Shaded}

\begin{verbatim}
## Linear mixed model fit by maximum likelihood  ['lmerMod']
## Formula: e.rel.130 ~ slurry.type + (1 | trial)
##    Data: isummul
## 
##       AIC       BIC    logLik -2*log(L)  df.resid 
##     -36.9     -30.7      22.5     -44.9        31 
## 
## Scaled residuals: 
##     Min      1Q  Median      3Q     Max 
## -1.6482 -0.7586 -0.2336  0.7202  2.2756 
## 
## Random effects:
##  Groups   Name        Variance Std.Dev.
##  trial    (Intercept) 0.00000  0.0000  
##  Residual             0.01623  0.1274  
## Number of obs: 35, groups:  trial, 6
## 
## Fixed effects:
##               Estimate Std. Error t value
## (Intercept)    0.29375    0.03090   9.508
## slurry.typePS -0.15064    0.04308  -3.497
## 
## Correlation of Fixed Effects:
##             (Intr)
## slrry.typPS -0.717
## optimizer (nloptwrap) convergence code: 0 (OK)
## boundary (singular) fit: see help('isSingular')
\end{verbatim}

\begin{Shaded}
\begin{Highlighting}[]
\FunctionTok{anova}\NormalTok{(mod3, mod3c, }\AttributeTok{test =} \StringTok{\textquotesingle{}Chisq\textquotesingle{}}\NormalTok{)}
\end{Highlighting}
\end{Shaded}

\begin{verbatim}
## Data: isummul
## Models:
## mod3c: e.rel.130 ~ slurry.type + (1 | trial)
## mod3: e.rel.130 ~ separation * slurry.type + (1 | trial)
##       npar     AIC     BIC logLik -2*log(L)  Chisq Df Pr(>Chisq)    
## mod3c    4 -36.912 -30.691 22.456   -44.912                         
## mod3     6 -76.226 -66.894 44.113   -88.226 43.314  2  3.931e-10 ***
## ---
## Signif. codes:  0 '***' 0.001 '**' 0.01 '*' 0.05 '.' 0.1 ' ' 1
\end{verbatim}

We still have the singular fit issue, but a very clear effect.

\section{Conclusions:}\label{conclusions}

\begin{itemize}
\tightlist
\item
  Separation affects emission
\item
  The reduction for the solid fraction differs between slurry types
\item
  The liquid fraction has lower emission than whole slurry for both, and
  the reduction is not clearly different between slurry types
\item
  For digestate, both liquid and solid fractions have lower emission
  than whole slurry, but are not clearly different from each other, if
  we consider that the separation effect varies among trials (and it
  clearly does, from mod2)
\item
  For pig slurry, whole slurry and the solid fraction are not clearly
  different, but liquid is clearly different (lower) than both
\end{itemize}

\end{document}
